\answer{ex:16:1}{Le sang est un circuit RC avec $\tau=t_{1/2}/\ln2=14.4$~min; rapport du maximum au minimum $2^{T/t_{1/2}}=64$; harmonique fondamentale atténuée à $0.55$; rapport $2$ pour $t_{1/2}=T=60$~min.}
\answer{ex:16:2}{$\varphi^*=\arcsin\bigl((\omega-\omega_0)/K_\varphi\bigr)\approx41$~min ($\omega-\omega_0\approx0.047$~rad/j), temps de relaxation $1/(K_\varphi\cos\varphi^*)\approx3.9$~j; après un décalage de $+6$ et de $-6$~h, $7.7$ et $7.9$~j.}
\answer{ex:16:3}{$K_c(n)=\sec^n(\pi/n)$: $8$ et $2.37$, erreur de $11\%$ et $30\%$; quand $n\to\infty$, $K_c\to1$ et l'erreur $\to50\%$.}
\answer{ex:16:4}{Les commandes extrêmes (accélérateur $80$, frein $0$) donnent $f=120$ en $0.76$~s; \og accélérateur seul\fg{}: $2.33$~s, trois fois plus.}
\answer{ex:16:5}{$3\arctan(\omega\tau)+\omega d=\pi$, $K_c=(1+\omega^2\tau^2)^{3/2}$: $K_c=8$, $3.54$, $2.50$, $1.76$ et périodes $3.63$, $5.46$, $6.86$, $9.27\,\tau$; à $0.9K_c$ les oscillations s'amortissent, à $1.1K_c$ elles s'installent.}
\answer{ex:16:6}{$0.60\bar c$, $0.87\bar c$, $0.90\bar c$, limite $0.91\bar c$; à concentration constante, la réponse est nulle: un tel destinataire ne lit que les bouffées.}
\answer{ex:16:7}{Question ouverte. Avec la rétroaction, le rythme n'existe que pour $K>8$; sa période est proportionnelle au $\tau$ des étages et dépend à peine de $K$ ($3.64\,\tau$ pour $K=8.5$, $4.23\,\tau$ pour $K=40$): changer $K$ d'un facteur une fois et demie décale la période de $2$ à $4\%$, moins que l'erreur de mesure. Ce qui tranche, c'est l'ouverture de la boucle (une hormone finale constante à l'entrée du premier étage supprime un tel rythme) ou la réponse à un bref ajout d'hormone à différentes phases du cycle.}
